% sections/bureau.tex: Section 5 (opening, 5.1 identification, 5.2 the payment margin). Numbers: results/panel_windows_auto.csv (T1l),
% identification_within_person.csv (I6), spec_scores.csv (bridge_spec_score.R), tab_benchmark (make_benchmark_table.R),
% runoff_model_rho_co_h0.22.csv, mc_runoff_model_grid_h0.22.csv, mc_persistent_*.csv.
\section{Estimates from the Credit-Bureau Panel}\label{sec:autodata}

Three margins and three samples recur through the rest of the paper. The \emph{payment margin} is the response of default to the monthly payment, every payment moved at once, with the balance moving with the payment at a fixed contract length: the elasticity of default within a window to the payment. The \emph{term margin} is the response of default to the contract's length at a fixed payment, which adds payments at the far end of the contract and, with them, a larger balance and more months in which a default can occur. The \emph{equity margin} is the response of mortgage default to the change in the house's value since origination. The three samples are the one-percent consumer panel, on which the population estimates run; the ten-percent trade archive, on which the within-borrower designs run; and the covariate-matched ten-percent extract, on which the specification search and the deficiency contrast run (Section \ref{sec:data}). Each of the designs the bureau supports has its own subsection below; the designs that do not identify the rate, and why, are recorded in Appendix \ref{appendix:strategies}.

\subsection{Identification}\label{sec:identification}

The population estimator is a linear probability model of default within $K$ months of origination on the logarithms of the monthly payment and the contract length,
\begin{equation}
\label{eq:main_reg}
\mathbf{1}[\text{default within } K]_{i} = \gamma_m \log m_i + \gamma_T \log T_i + \alpha_{\ell(i), y(i)} + \delta_{z(i), y(i)} + \eta_{v(i), q(i)} + \varepsilon_{i},
\end{equation}
with fixed effects for the lender by origination year, the three-digit zip code by origination year, and the credit-score ventile by income quartile, standard errors clustered by lender-year, and each coefficient divided by the window's default rate to give an elasticity. The identifying variation comes from two sources, and the fixed-effects regression is the vehicle that takes it to every loan. The first source is the published experiments and quasi-experiments of Section \ref{sec:quasi}, in which a rule moved a payment. The second is the within-borrower design: the same regression with borrower fixed effects, on borrowers who hold two or more auto loans in the ten-percent archive, in which the comparison is between two loans of the same person and nothing stable about her can explain the difference. The population regression is then what gives heterogeneity by score, income, place and cohort with power, and it is interpreted as causal under one assumption, stated next together with the sign and size of the bias if it fails and the evidence that the two identified sources and the population regression agree.

\paragraph{The assumption.} The payment coefficient is the response of Proposition \ref{prop:sensT} to the payment due now if, within a cell, $\log m$ is uncorrelated with the determinants of default that the cell does not absorb. The institutional argument for the assumption is that auto lenders set payment and term through underwriting rules that map the credit score, the vehicle and the down payment to a menu of allowable contracts, so that conditional on the score band and the local market the residual variation in $(m, T)$ reflects vehicle choice, dealer negotiation and the lender's menu rather than information about the borrower's future default that the score does not already contain.

\paragraph{The sign and the size of the bias.} If the assumption fails in the direction the institution suggests, a lender that observed default risk beyond the score would shorten terms and reduce payments for the riskier borrowers, which biases the payment coefficient upward and the term coefficient toward zero. The ratio of the term response to the payment response, the object the formula of equation \eqref{eq:window} maps to a rate, is then biased downward, and the rate upward: under that form of selection the population design overstates the rate. The within-borrower design removes every stable trait of the borrower and is the check on the size. Its payment elasticity is $0.44$ within twelve months and $0.53$ within twenty-four (Table \ref{tab:main_margins}, lower panel), against $0.45$ and $0.45$ for the population design on the one-percent panel and $0.37$ for the same design at twelve months on the covariate-matched extract: the population estimate equals the within-borrower one within sampling error at twelve months and is below it at twenty-four, so the gap, where there is one, has the opposite sign to the one lender screening on private information would produce, and it is small.

\paragraph{The menu of designs.} The cells of equation \eqref{eq:main_reg} were chosen among ten fixed-effects designs that span the nesting the data support, from origination-year effects alone to lender-by-zip3-by-year cells, by the design whose estimates come closest to the within-borrower estimates wherever the two can be compared, on 88 group-by-outcome comparisons on the covariate-matched extract (Appendix \ref{appendix:bridge}, Table \ref{tab:spec_search}). Across the ten designs the pooled twelve-month payment elasticity runs from $0.17$ with zip3-by-year cells alone to $0.41$ with linear controls for score, income and age; the coarse designs are low because within a place and year the borrowers who take larger payments are the ones with higher incomes and scores, who default less, so that without score and income cells the payment coefficient absorbs that sorting and is biased toward zero. The chosen design, with score-ventile-by-income-quartile cells added to lender-year and zip3-year cells, has the smallest mean gap to the within-borrower estimates on the term margin, $0.04$, and a gap of $0.12$ on the payment margin. The gap is not the same in every group: the test that it is constant across the nineteen groups where both designs exist rejects ($\chi^2$ of $1114.6$ on 166 degrees of freedom), so the population estimates by group are reported as they are, with the pooled gap stated, and the within-borrower estimate is the one the paper's rate rests on. Five further designs that hold the loan's balance fixed measure a different object, the response to the interest rate, and are reported with the benchmark table below. A coefficient chosen by its agreement with a benchmark is a post-selection estimate, and Appendix \ref{appendix:bridge} states the caveat for the standard errors of equation \eqref{eq:main_reg}; two instruments constructed from lenders' menus fail the exclusion restriction, and Appendix \ref{appendix:ivvalidation} reports them with the difference estimator that would have transported their information to the full sample.

\paragraph{The published benchmarks.} Table \ref{tab:benchmark} sets the payment responses of the published experiments and quasi-experiments beside the bureau's, in the units of an elasticity of default to an increase in the payment, at the two margins the literature distinguishes. Panel A is the paper's margin, a payment change with the balance moving, as when a larger loan or a longer contract is taken at the same rate. Panel B is the rate margin, a payment change at a fixed balance, as at an adjustable-rate reset or a modification that lowers the rate; there the bureau's comparable object is the within-borrower regression with the log balance added, and by the amortization identity that ties the payment, the term and the balance, its coefficients measure the response to the contract rate (Appendix \ref{appendix:ivvalidation}). On the rate margin the published elasticities of default to the payment run from $0.03$ and $0.3$ in two card studies through $1.2$ to $2.5$ for hybrid-mortgage resets, HAMP payment cuts and the subprime auto rate, to $5$ to $6.4$ for resets on bankruptcy, installment changes and renewals, and the bureau's within-borrower estimate with the balance held fixed, $2.01$ ($0.02$) at twelve months and $1.99$ ($0.01$) at twenty-four, sits inside that range. On the paper's own margin the published record is thin and does not agree with itself, $1.7$ ($0.1$) for a subprime auto loan-size hazard, $1.2$ ($0.3$) for a HAMP payment cut that extends the maturity and $-1.0$ ($0.5$) for the Dobbie--Song rescheduling, and the bureau's within-borrower $0.44$ ($0.005$) and $0.53$ ($0.003$) lie between them. Among the population designs the cells with a score-by-income dimension come closest to the within-borrower estimate, $0.41$ and $0.37$ against $0.44$, with the one-percent panel at $0.45$, while designs with place or lender cells alone give $0.17$ to $0.25$, half the within-borrower value; with the balance held fixed the population designs give $2.33$ to $2.42$ against $2.01$ within borrower, the excess being the cross-borrower selection in the rate a lender prices. Appendix \ref{appendix:experiments} states the conversion of every published row.


\subsection{The Payment Margin}\label{sec:mainresults}

Table \ref{tab:main_margins} reports both margins of equation \eqref{eq:main_reg} for exposure windows from six months to the life of the loan, and, below them, the same regression with borrower fixed effects on borrowers holding two or more auto loans; Appendix \ref{appendix:delta_method} states the inference for each object the paper reports with a standard error.

\begin{table}[H]
\centering
\caption{The two margins of default by exposure window}
\label{tab:main_margins}
\small
\begin{tabular}{lrrrrr}
\toprule
Default within & Loans & Defaults & Rate & Payment elasticity & Term elasticity \\
\midrule
\multicolumn{6}{l}{\emph{Charge-off}} \\
6 months & 7,496,292 & 6,211 & 0.08\% & 0.42 (0.04) & +0.11 (0.10) \\
12 months & 7,355,247 & 45,946 & 0.62\% & 0.45 (0.02) & +0.01 (0.05) \\
18 months & 7,163,011 & 89,466 & 1.25\% & 0.44 (0.02) & +0.05 (0.04) \\
24 months & 6,971,966 & 129,722 & 1.86\% & 0.45 (0.02) & +0.11 (0.04) \\
36 months & 6,416,292 & 181,962 & 2.84\% & 0.45 (0.02) & +0.18 (0.04) \\
48 months & 5,073,953 & 193,493 & 3.81\% & 0.46 (0.02) & +0.07 (0.04) \\
Life of the loan & 7,623,563 & 365,312 & 4.79\% & 0.35 (0.01) & +0.27 (0.02) \\
\multicolumn{6}{l}{\emph{First 90 days past due}} \\
6 months & 7,496,292 & 11,270 & 0.15\% & 0.39 (0.03) & -0.48 (0.08) \\
12 months & 7,355,247 & 53,651 & 0.73\% & 0.35 (0.02) & -0.27 (0.04) \\
18 months & 7,163,011 & 90,742 & 1.27\% & 0.35 (0.02) & -0.16 (0.04) \\
24 months & 6,971,966 & 121,568 & 1.74\% & 0.35 (0.02) & -0.10 (0.03) \\
36 months & 6,416,292 & 156,285 & 2.44\% & 0.35 (0.02) & -0.02 (0.03) \\
48 months & 5,073,953 & 152,911 & 3.01\% & 0.36 (0.02) & -0.17 (0.03) \\
Life of the loan & 7,623,563 & 285,139 & 3.74\% & 0.27 (0.01) & +0.21 (0.02) \\
\midrule
\multicolumn{6}{l}{\emph{Within borrower, charge-off (two or more auto loans; borrower and origination-year effects; ten-percent extract)}} \\
12 months & 62,519,528 & 423,757 & 0.68\% & 0.44 (0.005) & -0.31 (0.009) \\
24 months & 59,155,772 & 1,001,773 & 1.69\% & 0.53 (0.003) & +0.30 (0.005) \\
Life of the loan & 64,922,586 & 2,221,031 & 3.42\% & 0.46 (0.002) & +0.95 (0.003) \\
\bottomrule
\end{tabular}

\vspace{0.5em}
\begin{minipage}{0.92\textwidth}
\noindent \small \textit{Notes:} Elasticities of the default rate within the stated window with respect to the monthly payment and the contract length, from equation \eqref{eq:main_reg} on the auto panel, for charge-off and for the first ninety-day delinquency; standard errors in parentheses, clustered by lender-year. The lower panel absorbs borrower and origination-year effects on the multi-loan subsample of the ten-percent archive, the specification of every within-borrower estimate in the paper; the specification that absorbs lender-year in place of origination-year effects is reported in Appendix \ref{appendix:ivvalidation}.
\end{minipage}
% replication: Code/type/T1l_windows_panel.R, Code/identification/I6_within_person.R, Code/exhibits/make_main_table.R
\end{table}

\begin{table}[p]
\centering
\caption{Published payment responses beside the bureau's, by margin}
\label{tab:benchmark}
\footnotesize
% tables/tab_benchmark.tex — generated by Code/exhibits/make_benchmark_table.R from results/causal_ledger_literature.csv,
% identification_within_person.csv, identification_within_person_by_group_logL.csv, spec_search.csv, spec_search_v2.csv,
% panel_windows_auto.csv; do not edit by hand.
\setlength{\tabcolsep}{3pt}
\begin{tabularx}{\textwidth}{@{}>{\raggedright\arraybackslash}p{0.36\textwidth}>{\raggedright\arraybackslash}X>{\centering\arraybackslash}p{0.07\textwidth}>{\raggedleft\arraybackslash}p{0.13\textwidth}@{}}
\toprule
Study or design & Population and outcome & Horizon (months) & Elasticity (s.e.) \\
\midrule
\multicolumn{4}{@{}l}{\textbf{Panel A. Balance moves with the payment (the paper's margin)}} \\
\multicolumn{4}{@{}l}{\emph{Published}} \\
\citet{adams_liquidity_2009}, loan size & Subprime auto, one lender; default hazard & 40 & 1.71 (0.11) \\
\citet{ganong_liquidity_2020}, design 2 & HAMP payment cut; 90+ day delinquency & 24 & 1.18 (0.25) \\
\citet{dobbie_targeted_2020}, payment arm & Debt-management plans; 5-year bankruptcy & 60 & $-1.04$ (0.50) \\
\citet{dobbie_targeted_2020}, year 1 & Same experiment; bankruptcy in year 1 & 12 & $-0.73$ (0.58) \\
\citet{dobbie_targeted_2017}, median arms & Same experiment, 2017 arms; bankruptcy & 60 & $-1.06$ (0.64) \\
\addlinespace
\multicolumn{4}{@{}l}{\emph{Credit bureau, this paper}} \\
Within borrower, year cells & Auto, multi-loan borrowers; default & 12 & 0.44 (0.005) \\
 &  & 24 & 0.53 (0.003) \\
 &  & life & 0.46 (0.002) \\
Within borrower, lender-year &  & 12 & 0.51 (0.005) \\
 &  & 24 & 0.57 (0.003) \\
 &  & life & 0.47 (0.002) \\
S1: origination year & Auto, covariate-matched extract; default & 12 & 0.22 (0.03) \\
S2: zip3 $\times$ year &  & 12 & 0.17 (0.03) \\
S4: zip3 $\times$ year, lender $\times$ year &  & 12 & 0.24 (0.02) \\
S6: S4, score ventile $\times$ inc.\ quartile &  & 12 & 0.37 (0.02) \\
S8: S4, score, log income, age &  & 12 & 0.41 (0.02) \\
S9: lender $\times$ zip3 $\times$ year &  & 12 & 0.25 (0.02) \\
One-percent panel, S6 cells & Auto, one-percent panel; charge-off & 12 & 0.45 (0.02) \\
 &  & 24 & 0.45 (0.02) \\
\midrule
\multicolumn{4}{@{}l}{\textbf{Panel B. Balance fixed (the rate margin)}} \\
\multicolumn{4}{@{}l}{\emph{Published}} \\
\citet{adams_liquidity_2009}, APR & Subprime auto, one lender; default hazard & 40 & 1.47 (0.13) \\
\citet{fuster_payment_2017}, 3--3.5 pp cut & Alt-A interest-only ARMs; 60+ day hazard & 60 & 1.22 (0.10) \\
\citet{keys_mortgage_2014}, 5/1 vs 7/1 & Conforming hybrid ARMs; 60+ day & 24 & 2.47 (1.37) \\
\citet{indarte_moral_2023}, ARM resets & Indexed ARMs; bankruptcy in the reset year & 12 & 5.20 (1.47) \\
\citet{tracy_payment_2016}, payment cut & GSE prime ARMs, LTV above 80; 90+ day & 1 & 2.42 (0.16) \\
\citet{kartashova_how_2022}, renewals & Canadian 5-year mortgages; 60-day & 24 & 6.36 (1.16) \\
\citet{byrne_how_2022}, installments & Irish tracker mortgages; 90+ day hazard & 12 & 5.64 (0.28) \\
\citet{scharlemann_effect_2022}, step-up & HAMP modifications, step-up; default & 8 & 1.59 (0.38) \\
\citet{castellanos_contract_2026}, APR arm & Randomized card contracts, Mexico; default & 26 & 0.33 (0.07) \\
\citet{castellanos_contract_2026}, minimum & Same experiment; minimum 5 to 10\% & 26 & 0.31 (0.15) \\
\citet{mueller_increasing_2022}, IDR & Randomized IDR application; new 60+ day & 1 & 1.91 (n.r.) \\
\citet{dastous_liquidity_2017}, minimum & Card accounts, one bank; new delinquency & 12 & 0.03 (0.02) \\
\addlinespace
\multicolumn{4}{@{}l}{\emph{Credit bureau, this paper}} \\
Within borrower, year cells, log $L$ fixed & Auto, multi-loan borrowers; default & 12 & 2.01 (0.02) \\
 &  & 24 & 1.99 (0.01) \\
 &  & life & 1.93 (0.009) \\
S11: S6, log $L$ & Auto, covariate-matched extract; default & 12 & 2.40 (0.19) \\
S12: S4, ventile $\times$ inc.\ decile, log $L$ &  & 12 & 2.40 (0.19) \\
S13: zip3 $\times$ year, lender $\times$ tercile $\times$ year, ventile $\times$ decile, log $L$ &  & 12 & 2.33 (0.19) \\
S14: S4, decile $\times$ decile, log $L$, log inc., age &  & 12 & 2.42 (0.19) \\
\bottomrule
\end{tabularx}

\vspace{0.2em}
\begin{minipage}{0.95\textwidth}
\noindent \scriptsize \textit{Notes:} Elasticities of default to an increase in the payment, with standard errors in parentheses. Published rows are the studies' reported estimates converted to elasticities with the study's own base rate and payment change (Appendix \ref{appendix:experiments}). Panel A: the balance moves with the payment at a fixed term (the paper's margin). Panel B: the balance is fixed and the payment moves with the interest rate (the rate margin). Bureau rows: within borrower absorbs borrower and origination-year effects on the multi-loan borrowers of the ten-percent archive; the fixed-effects designs are the candidates of Appendix \ref{appendix:bridge} on the covariate-matched extract at twelve months; the one-percent panel row is equation \eqref{eq:main_reg}.
\end{minipage}
% replication: Code/exhibits/make_benchmark_table.R; results/causal_ledger_literature.csv, identification_within_person.csv, identification_within_person_by_group_logL.csv, spec_search.csv, spec_search_v2.csv, panel_windows_auto.csv
\end{table}

The table establishes two facts. The payment margin is a stable parameter. Its elasticity is $0.45$ within twelve months and $0.45$ within 24 on the population under the charge-off definition, $0.35$ under the ninety-day definition, and $0.44$ and $0.53$ within borrower; Section \ref{sec:heterogeneity} shows it is the same in every origination cohort from 2010 to 2024. The term margin is not one parameter. At a fixed payment a longer contract lowers early delinquency and raises lifetime default, the elasticity rising from $-0.27$ within twelve months to $+0.21$ over the life of the loan under the ninety-day definition and from $-0.31$ to $+0.95$ within borrower, and Section \ref{sec:termmargin} shows this to be the difference of a preference channel and a penalty that scales with the balance owed, plus the exposure that a longer contract adds.

\paragraph{The payment margin over the life of the loan.} Among loans that survive to each loan age, the elasticity of the default hazard to the payment declines as the contract shortens. Figure \ref{fig:runoff_profiles} in Appendix \ref{appendix:mcliquidity} shows it on 61- to 72-month contracts, pooled and by score tercile and income quartile: the pooled elasticity is $0.65$ in the first six months, $0.63$ through the second year, $0.54$ in the fourth and $0.49$ in the fifth, a ratio of the last bin to the first of $0.75$ under the charge-off definition and $0.67$ under the ninety-day definition, with the same shape in every group and a level that differs, $0.60$ on average in the bottom score tercile against $0.94$ in the middle. At loan age $s$ the incentive to default weighs the remaining payments, $W_s = 1 + \lambda \sum_{j=1}^{T-s} D(j)\,Q(j)$, against the balance that would be charged off, and a first-order formula in which $W_s$ declines at the rate $\rho + h$ would attribute a decline of $0.75$ over five years to an annual rate near $0.5$ (Appendix \ref{appendix:formula}). The stopping model of Section \ref{sec:model}, simulated with the penalty calibrated to the sample's default rate, produces a decline of $0.87$ to $0.99$ at every rate on a grid from zero to $1.2$ when income is independent across months, and of $0.22$ when income is persistent, because the borrower who continues keeps the option to default later and the far payments enter her decision only with the weight of the states in which she will still be paying them (Remark \ref{rem:option}); the shaded band in the figure is the independent-income range scaled to each group's first bin. Neither calibration reaches the level of the payment elasticity, $4.5$ in the model against $0.45$ in the data, and Appendix \ref{appendix:mcliquidity} reports the comparison by group and the two calibrations beside the data. The decline over the life of the loan is consistent with the forward-looking stopping problem the theory describes, and the paper does not use it as a measure of the rate in either direction; the rate comes from the within-borrower design of the next subsection.
